% ============================================================================
% APPENDIX MN: LIT-NIEDERLE - Gender, Competition & Market Design
% ============================================================================
% Category: Literature Integration (LIT-)
% Language: English
% EBF Integration: Gender differences, competition, and matching markets
% ============================================================================

\chapter*{Appendix MN: Muriel Niederle Research Integration}
\addcontentsline{toc}{chapter}{Appendix MN: LIT-NIEDERLE - Gender, Competition \& Market Design}
\label{app:niederle}

% ============================================================================
% HEADER BLOCK
% ============================================================================
\begin{tcolorbox}[colback=white,colframe=black!75!white,title=Appendix Metadata]
\textbf{Appendix:} MN -- LIT-NIEDERLE \\
\textbf{Category:} Literature Integration (LIT-) \\
\textbf{Version:} 1.0 (January 2026) \\
\textbf{Dependencies:} LIT-GNEEZY, LIT-ROTH, CONTEXT-DEMOGRAPHIC, DOMAIN-HR \\
\textbf{Status:} Complete
\end{tcolorbox}

% ============================================================================
% CROSS-REFERENCE MAP
% ============================================================================
\begin{tcolorbox}[colback=blue!5!white,colframe=blue!75!black,title=Cross-Reference Map]
\textbf{Builds On:}
\begin{itemize}
    \item Appendix GN (LIT-GNEEZY): Gender and competition foundations
    \item Appendix RO (LIT-ROTH): Market design principles
    \item Appendix SA (LIT-AMBÜHL): Organ market research collaboration
    \item Appendix AI (CONTEXT-DEMOGRAPHIC): Demographic context effects
\end{itemize}

\textbf{Extends To:}
\begin{itemize}
    \item Appendix AG (DOMAIN-HR): Workplace gender dynamics
    \item Appendix AJ (DOMAIN-EDUCATION): Educational competition
    \item Appendix S (PREDICT-POLICY): Policy interventions for gender equity
\end{itemize}

\textbf{Methods Connection:}
\begin{itemize}
    \item Appendix E (METHOD-E): Laboratory experiments
    \item Appendix AN (METHOD-LLMMC): Parameter estimation
\end{itemize}
\end{tcolorbox}

% ============================================================================
% CHAPTER LINKAGE
% ============================================================================
\begin{tcolorbox}[colback=green!5!white,colframe=green!75!black,title=Chapter Linkage]
\textbf{Primary:} Chapter 9 (Competition \& Incentives), Chapter 14 (HR Applications) \\
\textbf{Secondary:} Chapter 5 (Social Preferences), Chapter 10 (Context Effects)
\end{tcolorbox}

% ============================================================================
% ABSTRACT
% ============================================================================
\begin{abstract}
This appendix integrates Muriel Niederle's research on gender differences in competition and market design into the Evidence-Based Framework. Niederle's work (often with Lise Vesterlund) established that women systematically ``shy away'' from competition even when equally capable, contributing substantially to gender gaps in high-stakes environments. Her research demonstrates that this gap is malleable through institutional design, including affirmative action policies. In collaboration with Alvin Roth, she has also advanced market design, particularly for matching markets. This appendix provides critical parameters for CONTEXT-DEMOGRAPHIC and establishes evidence-based interventions for gender equity.
\end{abstract}

% ============================================================================
% QUICK REFERENCE
% ============================================================================
\begin{tcolorbox}[colback=gray!10!white,colframe=gray!50!black,title=Quick Reference]
Key terms defined in this appendix (see also Appendix G for complete Glossary):
\begin{description}
    \item[Competition Gap] Gender difference in willingness to enter competitive environments
    \item[Tournament Selection] Choosing competitive over piece-rate compensation
    \item[Affirmative Action] Policies to increase underrepresented group participation
    \item[Matching Market] Two-sided market where participants must be matched
\end{description}
\end{tcolorbox}

% ============================================================================
% FUNDAMENTAL QUESTION
% ============================================================================
% -----------------------------------------------------------------------------
% APPENDIX SCOPE BOX
% -----------------------------------------------------------------------------

\begin{tcolorbox}[
    colback=orange!5!white,
    colframe=orange!75!black,
    title={\textbf{Appendix Scope: Ziel / In-Scope / Out-of-Scope / Constraints}},
    fonttitle=\bfseries
]

\textbf{Ziel (Objective):}
\begin{itemize}[nosep]
    \item Why do equally capable women and men make different choices about competition, and how can institutions be designed to reduce inefficient gender gaps?
\end{itemize}

\textbf{In-Scope:}
\begin{itemize}[nosep]
    \item Fundamental Question
    \item Theoretical Framework
    \item Core Axioms
    \item Key Empirical Results
    \item EBF Integration
\end{itemize}

\textbf{Out-of-Scope (delegiert an andere Appendices):}
\begin{itemize}[nosep]
    \item LIT-GNEEZY $\rightarrow$ Appendix GN
    \item LIT-ROTH $\rightarrow$ Appendix RO
    \item LIT-AMBÜHL $\rightarrow$ Appendix SA
    \item CONTEXT-DEMOGRAPHIC $\rightarrow$ Appendix AI
    \item DOMAIN-HR $\rightarrow$ Appendix AG
\end{itemize}

\textbf{Constraints (Anwendungsgrenzen):}
\begin{itemize}[nosep]
    \item [TODO: Define when this appendix does NOT apply]
    \item [TODO: Define required prerequisites]
    \item [TODO: Define known limitations]
\end{itemize}

\textbf{Lieferobjekte (Deliverables):}
\begin{itemize}[nosep]
    \item 5 Table(s)
    \item 10 Equation(s)
    \item 6 Axiom(s)
    \item Worked Example(s)
\end{itemize}

\end{tcolorbox}


\section{Fundamental Question}

\begin{tcolorbox}[colback=yellow!10!white,colframe=orange!75!black,title=Central Research Question]
\textbf{Why do equally capable women and men make different choices about competition, and how can institutions be designed to reduce inefficient gender gaps?}
\end{tcolorbox}

This question addresses CONTEXT-DEMOGRAPHIC: How does gender interact with competitive environments?

% ============================================================================
% THEORETICAL FRAMEWORK
% ============================================================================
\section{Theoretical Framework}

\subsection{The Competition Selection Model}

Niederle's framework decomposes the gender gap in competitive entry:

\begin{equation}
P(\text{tournament})_i = f(\text{Ability}_i, \text{Risk}_i, \text{Confidence}_i, \text{Gender}_i)
\end{equation}

Where the gender effect operates through multiple channels:
\begin{itemize}
    \item Risk preferences (small effect)
    \item Overconfidence (moderate effect)
    \item Competition aversion (large effect)
\end{itemize}

\subsection{Decomposition of the Gender Gap}

\begin{equation}
\text{Gap} = \underbrace{\beta_1 \cdot \Delta\text{Ability}}_{0\%} + \underbrace{\beta_2 \cdot \Delta\text{Risk}}_{15\%} + \underbrace{\beta_3 \cdot \Delta\text{Confidence}}_{25\%} + \underbrace{\beta_4 \cdot \text{Competition Aversion}}_{60\%}
\end{equation}

Key insight: Most of the gap is not explained by ability or risk---it's competition-specific.

\subsection{Market Design Principles}

From Niederle-Roth collaboration:

\begin{equation}
\text{Market Performance} = f(\text{Thickness}, \text{Congestion}, \text{Safety})
\end{equation}

Where:
\begin{itemize}
    \item Thickness: Enough participants on both sides
    \item Congestion: Manageable number of offers/decisions
    \item Safety: Incentive to reveal true preferences
\end{itemize}

% ============================================================================
% AXIOMS
% ============================================================================
\section{Core Axioms}

\begin{tcolorbox}[colback=red!5!white,colframe=red!75!black,title=Axiom MN-1: Gender Competition Gap]
\textbf{Women are significantly less likely to select into competitive environments than equally-capable men.}

\begin{equation}
P(\text{tournament} | \text{Female, Top Quartile}) \approx 0.35 \quad \text{vs.} \quad P(\text{tournament} | \text{Male, Top Quartile}) \approx 0.75
\end{equation}

The 40 percentage point gap persists even among the highest performers.

\textit{Evidence:} Niederle \& Vesterlund (2007), QJE; replicated across 10+ studies.
\end{tcolorbox}

\begin{tcolorbox}[colback=red!5!white,colframe=red!75!black,title=Axiom MN-2: Overconfidence Channel]
\textbf{Men are more overconfident about relative performance, explaining part of the competition gap.}

\begin{equation}
\hat{P}(\text{win} | \text{Male}) - P(\text{win} | \text{Male}) > \hat{P}(\text{win} | \text{Female}) - P(\text{win} | \text{Female})
\end{equation}

Men overestimate their rank by $\sim$15 percentile points more than women.

\textit{Evidence:} Niederle \& Vesterlund (2007), belief elicitation data.
\end{tcolorbox}

\begin{tcolorbox}[colback=red!5!white,colframe=red!75!black,title=Axiom MN-3: Affirmative Action Efficacy]
\textbf{Affirmative action in competitive settings increases female entry without reducing average quality.}

\begin{equation}
\text{Quality}_{AA} \geq \text{Quality}_{No-AA} \quad \text{despite} \quad \text{Female Entry}_{AA} >> \text{Female Entry}_{No-AA}
\end{equation}

The marginal female entrant is often higher ability than the marginal male she replaces.

\textit{Evidence:} Niederle, Segal \& Vesterlund (2013), Management Science.
\end{tcolorbox}

\begin{tcolorbox}[colback=red!5!white,colframe=red!75!black,title=Axiom MN-4: Competition Aversion Independence]
\textbf{Competition aversion is distinct from risk aversion and largely explains the residual gender gap.}

\begin{equation}
\text{Gap} | \text{Control for Risk} \approx 0.30 \quad (\text{still significant})
\end{equation}

Women dislike competition \textit{per se}, not just the uncertainty of outcomes.

\textit{Evidence:} Niederle \& Vesterlund (2011), Annual Review of Economics.
\end{tcolorbox}

\begin{tcolorbox}[colback=red!5!white,colframe=red!75!black,title=Axiom MN-5: Market Unraveling]
\textbf{Without centralized clearinghouses, markets unravel through exploding offers and early contracting.}

\begin{equation}
t_{contract} \xrightarrow{\text{unraveling}} 0 \quad \Rightarrow \quad \text{Match Quality} \downarrow
\end{equation}

Early offers reduce information available for matching, harming efficiency.

\textit{Evidence:} Niederle \& Roth (2009), AEJ: Microeconomics.
\end{tcolorbox}

\begin{tcolorbox}[colback=red!5!white,colframe=red!75!black,title=Axiom MN-6: Institutional Malleability]
\textbf{The gender competition gap can be substantially reduced through institutional design.}

Effective interventions:
\begin{itemize}
    \item Quotas: Reduce gap by $\sim$50\%
    \item Team competition: Reduce gap by $\sim$30\%
    \item Feedback: Reduce gap by $\sim$20\%
\end{itemize}

\textit{Evidence:} Multiple studies synthesized in Niederle (2016).
\end{tcolorbox}

% ============================================================================
% KEY EMPIRICAL RESULTS
% ============================================================================
\section{Key Empirical Results}

\subsection{The Core Competition Experiment}

\begin{table}[h]
\centering
\begin{tabular}{lccc}
\hline
\textbf{Group} & \textbf{Chose Tournament} & \textbf{Performance (if competed)} & \textbf{Win Rate} \\
\hline
Men & 73\% & 11.2 problems & 45\% \\
Women & 35\% & 12.1 problems & 55\% \\
\hline
\end{tabular}
\caption{Competition entry and performance by gender}
\end{table}

\textbf{Key insight:} Women who enter competitions perform \textit{better} than men who enter.

\subsection{Affirmative Action Results}

\begin{table}[h]
\centering
\begin{tabular}{lccc}
\hline
\textbf{Condition} & \textbf{Female Entry} & \textbf{Avg Quality} & \textbf{Efficiency Loss} \\
\hline
No intervention & 35\% & 10.8 & -- \\
Quota (2 of 4) & 68\% & 10.9 & None \\
Quota (3 of 4) & 75\% & 10.4 & Small \\
\hline
\end{tabular}
\caption{Effects of affirmative action on entry and quality}
\end{table}

\subsection{Decomposition of Gap Sources}

\begin{table}[h]
\centering
\begin{tabular}{lcc}
\hline
\textbf{Factor} & \textbf{Share of Gap} & \textbf{Malleable?} \\
\hline
Ability differences & 0\% & N/A \\
Risk preferences & 15\% & Somewhat \\
Overconfidence & 25\% & Yes (feedback) \\
Competition aversion & 60\% & Yes (institutions) \\
\hline
\end{tabular}
\caption{Decomposition of gender competition gap}
\end{table}

% ============================================================================
% EBF INTEGRATION
% ============================================================================
\section{EBF Integration}

\subsection{CONTEXT-DEMOGRAPHIC Parameters}

Niederle's research provides critical parameters:

\begin{equation}
\Psi_{gender} = \begin{pmatrix}
\gamma_{comp} & \text{(competition aversion)} \\
\delta_{conf} & \text{(confidence gap)} \\
\eta_{AA} & \text{(affirmative action effect)} \\
\end{pmatrix}
\end{equation}

\begin{table}[h]
\centering
\begin{tabular}{lccl}
\hline
\textbf{Parameter} & \textbf{Estimate} & \textbf{SE} & \textbf{Source} \\
\hline
$\gamma_{comp}$ (gap) & 0.38 & 0.05 & QJE 2007 \\
$\delta_{conf}$ (overconfidence) & 0.15 & 0.04 & Multiple studies \\
$\eta_{AA}$ (quota effect) & 0.33 & 0.08 & Mgmt Sci 2013 \\
Institutional reduction & 0.50 & 0.10 & Meta-analysis \\
\hline
\end{tabular}
\caption{EBF parameters from Niederle research}
\end{table}

\subsection{10C Framework Mapping}

\begin{table}[h]
\centering
\begin{tabular}{lll}
\hline
\textbf{CORE} & \textbf{Niederle Contribution} & \textbf{Key Paper} \\
\hline
WHO (AAA) & Gender as welfare-relevant category & All work \\
WHAT (C) & Competition utility component & QJE 2007 \\
HOW (B) & Gender-competition interaction & Ann Rev 2011 \\
WHEN (V) & Institutional context effects & Mgmt Sci 2013 \\
WHERE (BBB) & Gap parameters & Multiple \\
AWARE (AU) & Overconfidence asymmetry & QJE 2007 \\
\hline
\end{tabular}
\caption{10C Framework mapping for Niederle research}
\end{table}

% ============================================================================
% WORKED EXAMPLE
% ============================================================================
\section{Worked Example: Designing Gender-Inclusive Hiring}

\begin{tcolorbox}[colback=blue!5!white,colframe=blue!75!black,title=Application: Tech Company Leadership Pipeline]
\textbf{Context:} A tech company has 40\% women in entry roles but only 15\% in leadership.

\textbf{Step 1: Diagnose Competition Points}

Identify where competitive selection occurs:
\begin{itemize}
    \item Promotion applications (self-nomination)
    \item Leadership development programs (competitive entry)
    \item Stretch assignments (volunteer basis)
\end{itemize}

\textbf{Step 2: Apply Niederle Parameters}

Expected female self-nomination rate:
\begin{equation}
P(\text{apply} | F) = P(\text{apply} | M) \times (1 - 0.38) = 0.62 \times P(\text{apply} | M)
\end{equation}

\textbf{Step 3: Design Interventions}

\begin{enumerate}
    \item \textbf{Nomination by others:} Remove self-selection (+20\% female candidates)
    \item \textbf{Cohort-based programs:} Team vs. individual competition (+15\%)
    \item \textbf{Transparent criteria:} Reduce ambiguity that favors overconfident (+10\%)
    \item \textbf{Modest quotas:} Ensure representation without quality loss
\end{enumerate}

\textbf{Step 4: Predict Outcomes}

With combined interventions:
\begin{align}
\text{Female Leadership (baseline)} &= 15\% \\
\text{Female Leadership (interventions)} &\approx 15\% \times 1.5 = 22.5\%
\end{align}

Over 5 years with pipeline effects: Target 30\%+ achievable.
\end{tcolorbox}

% ============================================================================
% IMPLICATIONS
% ============================================================================
\section{Implications for Practice}

\subsection{For Organizations}

\begin{enumerate}
    \item \textbf{Audit Competition Points:} Identify where self-selection reduces diversity
    \item \textbf{Nomination Systems:} Supplement self-nomination with peer/manager nomination
    \item \textbf{Team-Based Competition:} Reduce individual tournament structures
    \item \textbf{Feedback Loops:} Provide accurate performance information to reduce confidence gap
\end{enumerate}

\subsection{For Policy}

\begin{enumerate}
    \item \textbf{Quotas Work:} Properly designed quotas increase diversity without quality loss
    \item \textbf{Early Intervention:} Address competition aversion before high-stakes decisions
    \item \textbf{Market Design:} Centralized matching reduces discrimination opportunities
\end{enumerate}

% ============================================================================
% SUMMARY
% ============================================================================
\section{Summary}

\begin{tcolorbox}[colback=gray!10!white,colframe=gray!50!black,title=Key Takeaways]
\begin{enumerate}
    \item \textbf{38pp Gap:} Women 38 percentage points less likely to enter competition
    \item \textbf{Not Ability:} Gap exists even among top performers
    \item \textbf{60\% Competition Aversion:} Most of gap is competition-specific, not risk
    \item \textbf{Affirmative Action Works:} Quotas increase entry without reducing quality
    \item \textbf{Institutions Matter:} Gap is malleable through design (50\% reduction possible)
    \item \textbf{Market Design:} Centralized clearinghouses improve matching efficiency
\end{enumerate}

\textbf{Core Insight:} The gender gap in competitive environments is real but fixable through institutional design.
\end{tcolorbox}

% ============================================================================
% GLOSSARY
% ============================================================================
\section{Glossary}

Key terms defined in this appendix (see also Appendix G for complete Glossary):

\begin{description}
    \item[Competition Gap] Gender difference in tournament entry rates
    \item[Tournament Selection] Choosing winner-take-all over piece-rate pay
    \item[Competition Aversion] Disutility from competing, distinct from risk aversion
    \item[Affirmative Action] Reserved slots or preferential treatment for underrepresented groups
    \item[Matching Market] Market where both sides must agree (jobs, schools, organs)
    \item[Market Unraveling] Premature contracting that reduces match quality
    \item[Exploding Offer] Time-limited offer designed to prevent comparison shopping
\end{description}

% ============================================================================
% CRITICAL FOUNDATIONS
% ============================================================================
\section{Critical Foundations}

\begin{enumerate}
    \item \textbf{Gneezy, Niederle \& Rustichini (2003):} Original gender-competition finding
    \item \textbf{Croson \& Gneezy (2009):} Gender differences in preferences review
    \item \textbf{Roth \& Sotomayor (1990):} Two-sided matching theory
    \item \textbf{Bertrand (2011):} New perspectives on gender
    \item \textbf{Goldin (2014):} Grand gender convergence
\end{enumerate}

% ============================================================================
% OPEN ISSUES
% ============================================================================
\section{Open Issues}

\begin{enumerate}
    \item \textbf{Origins:} Nature vs. nurture in competition aversion?
    \item \textbf{Age Effects:} Does the gap change over the career?
    \item \textbf{Domain Specificity:} Is the gap larger in male-typed domains?
    \item \textbf{Long-term AA Effects:} Do quotas change underlying preferences?
    \item \textbf{Intersectionality:} How does race interact with gender in competition?
\end{enumerate}

% ============================================================================
% REFERENCES
% ============================================================================
\section{References}

See Master Bibliography (bcm2\_master\_references.bib) for complete citations.

\nocite{bcm_master}

\textbf{Primary Niederle Sources} (23 papers integrated):

\begin{itemize}
    \item Niederle, M. \& Vesterlund, L. (2007). Do Women Shy Away from Competition? \textit{QJE}, 122(3).
    \item Niederle, M. \& Vesterlund, L. (2011). Gender and Competition. \textit{Annual Review of Economics}, 3.
    \item Niederle, M., Segal, C. \& Vesterlund, L. (2013). How Costly Is Diversity? \textit{Management Science}, 59(1).
    \item Niederle, M. \& Vesterlund, L. (2010). Explaining the Gender Gap in Math. \textit{JEP}, 24(2).
    \item Gneezy, U., Leonard, K. \& List, J. (2009). Gender Differences: Matrilineal vs Patriarchal. \textit{Econometrica}.
    \item Niederle, M. \& Roth, A. (2009). Market Culture: Exploding Offers. \textit{AEJ: Micro}, 1(2).
    \item Ambühl, S., Niederle, M. \& Roth, A. (2023). Organ Markets and Beliefs. \textit{QJE}, 138(1).
\end{itemize}

\end{document}
